\section{Representable functors and the Yoneda lemma}
\source{fga-explained:page-0023}\source{fga-explained:page-0024}\source{fga-explained:page-0025}\source{fga-explained:page-0026}\source{fga-explained:page-0027}\source{fga-explained:page-0028}\source{fga-explained:page-0029}\source{fga-explained:page-0030}\source{fga-explained:page-0031}\source{fga-explained:page-0032}\source{fga-explained:page-0033}\source{fga-explained:page-0034}\source{fga-explained:page-0035}\source{fga-explained:page-0036}\source{fga-explained:page-0037}\source{fga-explained:page-0038}\source{fga-explained:page-0039}\source{fga-explained:page-0040}

For a category $\mathcal C$ and $X\in\mathcal C$, put
$h_X(U)=\Hom_{\mathcal C}(U,X)$.  A morphism $X\to Y$ induces a natural
transformation $h_X\to h_Y$ by composition.

\begin{definition}[representable functor]\label{fga-def-2-1}\dcref{2.1}\source{fga-explained:page-0023}\uses{}
A contravariant set-valued functor $F$ on $\mathcal C$ is representable if it
is isomorphic to $h_X$ for some $X$; then $X$ is a representing object.
\end{definition}\begin{theorem}[Yoneda]\label{fga-thm-2-1}\source{fga-explained:page-0024}\uses{fga-def-2-1}
For every $X\in\mathcal C$ and contravariant functor $F$ there is a natural
bijection $\operatorname{Nat}(h_X,F)\simeq F(X)$.  A transformation is sent
to the image of $\operatorname{id}_X$, and $a\in F(X)$ gives the transformation
$f:U\to X\mapsto F(f)(a)$.  In particular
$\Hom_{\mathcal C}(X,Y)\simeq\operatorname{Nat}(h_X,h_Y)$.
\end{theorem}
\begin{proof}
The two constructions are inverse by naturality: evaluating the transformation
associated to $a$ at $\operatorname{id}_X$ returns $a$, while naturality at
$f:U\to X$ recovers its value on $f$.
\end{proof}

\begin{definition}[universal object]\label{fga-def-2-2}\dcref{2.2}\source{fga-explained:page-0025}\uses{fga-thm-2-1}
A universal object for $F$ is a pair $(X,\xi)$ with $\xi\in F(X)$ such that
for every $U$ and $a\in F(U)$ there is a unique $f:U\to X$ with
$F(f)(\xi)=a$.
\end{definition}
\begin{proposition}\label{fga-prop-2-3}\dcref{2.3}\source{fga-explained:page-0025}\uses{fga-def-2-2,fga-thm-2-1}
$F$ is representable if and only if it has a universal object.
\end{proposition}
\begin{proof}
The element corresponding to an isomorphism $h_X\to F$ is universal.  A
universal element defines the Yoneda isomorphism $h_X\to F$.
\end{proof}

\begin{remark}[uniqueness of representing objects]\label{fga-rem-2-4}\source{fga-explained:page-0025}\uses{fga-thm-2-1}
Representing objects are unique up to a unique isomorphism, and the functor
$X\mapsto h_X$ is fully faithful.
\end{remark}

The functors of functions, invertible functions, generated invertible quotients,
relative projective space, and locally free rank-$r$ quotients are represented
respectively by affine space, $\mathbb G_m$, projective space, $\Proj\operatorname{Sym}
\mathcal M$, and a relative Grassmannian.  These examples fix the functor-of-
points interpretation used throughout the volume.

\section{Group objects}
\source{fga-explained:page-0028}\source{fga-explained:page-0029}\source{fga-explained:page-0030}\source{fga-explained:page-0031}\source{fga-explained:page-0032}\source{fga-explained:page-0033}\source{fga-explained:page-0034}

\begin{definition}[group object]\label{fga-def-2-11}\dcref{2.11}\source{fga-explained:page-0042}\uses{}
In a category with finite products, a group object $G$ consists of multiplication
$m:G\times G\to G$, identity $e:1\to G$, and inverse $i:G\to G$ satisfying
the usual associative, unit, and inverse diagrams.
\end{definition}\begin{proposition}\label{fga-prop-2-12}\dcref{2.12}\source{fga-explained:page-0042}\uses{fga-def-2-11,fga-thm-2-1}
For every object $U$, the set $G(U)$ is a group, functorially in $U$; conversely
a representable group-valued functor gives a unique group-object structure.
\end{proposition}
\begin{proof}
Evaluate the structure diagrams on $U$ and use Yoneda for the converse.
\end{proof}

\begin{definition}[group-object homomorphism]\label{fga-def-2-13}\dcref{2.13}\source{fga-explained:page-0042}\uses{fga-def-2-11}
A morphism of group objects is a map commuting with multiplication, identity,
and inverse (equivalently, it induces group homomorphisms on all functor-of-
points groups).
\end{definition}

\begin{definition}[action]\label{fga-def-2-15}\dcref{2.15}\source{fga-explained:page-0044}\uses{fga-def-2-11}
A left action of a group-valued functor $G$ on a set-valued functor $F$ is a
natural transformation $G\times F\to F$ whose value at each $U$ is a group
action.  An action of a group object on $X$ means an action of $h_G$ on $h_X$.
\end{definition}
\begin{proposition}\label{fga-prop-2-16}\dcref{2.16}\source{fga-explained:page-0044}\uses{fga-def-2-15,fga-thm-2-1}
An action of $G$ on $X$ is equivalently a map $a:G\times X\to X$ satisfying
$a(e,x)=x$ and $a(g_1,a(g_2,x))=a(m(g_1,g_2),x)$.
\end{proposition}
\begin{proof}
The two displayed action diagrams become the ordinary action axioms after
evaluation on every test object; Yoneda then identifies the natural maps.
\end{proof}

\begin{definition}\label{fga-def-2-17}\dcref{2.17}\source{fga-explained:page-0045}\uses{fga-def-2-15}
An arrow between $G$-objects is $G$-equivariant when it commutes with the two
actions.
\end{definition}
\begin{proposition}\label{fga-prop-2-18}\dcref{2.18}\source{fga-explained:page-0046}\uses{fga-def-2-15,fga-thm-2-1}
Giving an action of a group-valued functor $G$ on $X$ is equivalent to giving a
natural transformation $G\to\Aut(X)$, where $\Aut(X)(U)$ is the group of
automorphisms of $U\times X$ over $U$.
\end{proposition}

\begin{definition}[discrete objects]\label{fga-def-2-19}\dcref{2.19}\source{fga-explained:page-0047}\uses{}
A category has discrete objects when it has a terminal object and all coproducts
of copies of it.  The coproduct indexed by a set is the associated discrete
object.
\end{definition}\begin{definition}[discrete group objects]\label{fga-def-2-20}\dcref{2.20}\source{fga-explained:page-0047}\uses{fga-def-2-19}
A category has discrete group objects when it has finite products, coproducts
$I\times U=\coprod_I U$, and the canonical map $I\times U\to A(I)\times U$
is an isomorphism for every set $I$ and object $U$.
\end{definition}
\begin{proposition}\label{fga-prop-2-21}\dcref{2.21}\source{fga-explained:page-0048}\uses{fga-def-2-20}
Under these hypotheses the discrete-object functor preserves finite products and
therefore sends groups to discrete group objects.
\end{proposition}
\begin{proposition}\label{fga-prop-2-22}\dcref{2.22}\source{fga-explained:page-0050}\uses{fga-prop-2-21,fga-def-2-15}
For a category with discrete group objects, an action of a discrete group on an
object is equivalent to an action of the associated discrete group object.
\end{proposition}

\section{Sheaves in Grothendieck topologies}
\source{fga-explained:page-0035}\source{fga-explained:page-0036}\source{fga-explained:page-0037}\source{fga-explained:page-0038}\source{fga-explained:page-0039}\source{fga-explained:page-0040}\source{fga-explained:page-0041}\source{fga-explained:page-0042}\source{fga-explained:page-0043}\source{fga-explained:page-0044}\source{fga-explained:page-0045}\source{fga-explained:page-0046}\source{fga-explained:page-0047}\source{fga-explained:page-0048}\source{fga-explained:page-0049}\source{fga-explained:page-0050}

\begin{definition}[Grothendieck topology]\label{fga-def-2-24}\dcref{2.24}\source{fga-explained:page-0075}\uses{}
A Grothendieck topology on $\mathcal C$ assigns to every object $U$ a class of
covering families $\{U_i\to U\}$ containing isomorphisms, stable under base
change, and stable under composition of coverings.  A category with such data is
a site.
\end{definition}The Zariski, etale, fppf and fpqc coverings of schemes are respectively jointly
surjective families of open immersions, etale maps locally of finite
presentation, flat maps locally of finite presentation, and families whose
coproduct map is faithfully flat and quasi-compact.

\begin{proposition}\label{fga-prop-2-33}\dcref{2.33}\source{fga-explained:page-0077}\uses{fga-def-1-3}
For a surjective morphism $f:X\to Y$, the following are equivalent: every
quasi-compact open of $Y$ is the image of a quasi-compact open of $X$; this
holds on an affine-open cover of $Y$; every $x\in X$ has a quasi-compact open
neighborhood mapping to an open neighborhood of $f(x)$; and every $x$ has such
a neighborhood with affine image.
\end{proposition}
\begin{definition}[fpqc morphism]\label{fga-def-2-34}\dcref{2.34}\source{fga-explained:page-0078}\uses{fga-prop-2-33,fga-def-1-7}
A faithfully flat morphism satisfying the equivalent conditions of
\cref{fga-prop-2-33} is fpqc.  An fpqc covering is a family whose coproduct map
is fpqc.
\end{definition}
\begin{proposition}\label{fga-prop-2-35}\dcref{2.35}\source{fga-explained:page-0078}\uses{fga-def-2-34}
Fpqc morphisms are stable under composition and base change, are local on the
target, include open faithfully flat maps and faithfully flat locally finitely
presented maps, and detect openness of subsets.
\end{proposition}
\begin{proposition}\label{fga-prop-2-36}\dcref{2.36}\source{fga-explained:page-0079}\uses{fga-prop-2-35}
Separatedness, quasi-compactness, finite presentation, properness, affineness,
finiteness, flatness, smoothness, unramifiedness, etaleness, embeddings and
closed embeddings descend along fpqc coverings.
\end{proposition}

\begin{definition}[sheaf]\label{fga-def-2-37}\dcref{2.37}\source{fga-explained:page-0080}\uses{fga-def-2-24}
A presheaf $F$ on a site is separated if sections agreeing after pullback to a
cover agree.  It is a sheaf if compatible sections $a_i\in F(U_i)$ on a
covering family glue to a unique $a\in F(U)$; compatibility means equality of
the two pullbacks to every $U_i\times_UU_j$.
\end{definition}
\begin{definition}[sieve]\label{fga-def-2-38}\dcref{2.38}\source{fga-explained:page-0080}\uses{fga-thm-2-1}
A sieve on $U$ is a subfunctor $S\subseteq h_U$ closed under precomposition.
The sieve generated by a family $\{U_i\to U\}$ consists of arrows factoring
through some $U_i$.
\end{definition}
\begin{proposition}\label{fga-prop-2-39}\dcref{2.39}\source{fga-explained:page-0081}\uses{fga-def-2-38,fga-thm-2-1}
For a covering family there is a canonical bijection
$\operatorname{Nat}(h_{\{U_i\}},F)\simeq F(\{U_i\})$, identifying compatible
families of sections with maps from the covering sieve.
\end{proposition}
\begin{corollary}\label{fga-cor-2-40}\dcref{2.40}\source{fga-explained:page-0081}\uses{fga-def-2-37,fga-prop-2-39}
$F$ is a sheaf exactly when $F(U)\to\operatorname{Nat}(h_U,F)$ factors through
every covering sieve by a bijection; it is separated exactly when these maps are
injective.
\end{corollary}

\begin{definition}\label{fga-def-2-41}\dcref{2.41}\source{fga-explained:page-0082}\uses{fga-def-2-24,fga-def-2-38}
A sieve belongs to a topology if it contains a generated covering sieve.
\end{definition}
\begin{proposition}\label{fga-prop-2-42}\dcref{2.42}\source{fga-explained:page-0082}\uses{fga-cor-2-40,fga-def-2-41}
The sheaf condition depends only on the associated Grothendieck topology and
not on the chosen pretopology.
\end{proposition}
\begin{proposition}\label{fga-prop-2-44}\dcref{2.44}\source{fga-explained:page-0083}\uses{fga-def-2-41}
Pullback of a covering sieve along any arrow is a covering sieve, and coverings
are stable under refinement.
\end{proposition}
\begin{definition}\label{fga-def-2-45}\dcref{2.45}\source{fga-explained:page-0084}\uses{fga-def-2-38}
The pullback of a family $\{U_i\to U\}$ along $V\to U$ is the family
$\{U_i\times_UV\to V\}$; two families are equivalent when they generate the
same sieve.
\end{definition}
\begin{proposition}\label{fga-prop-2-46}\dcref{2.46}\source{fga-explained:page-0084}\uses{fga-def-2-45,fga-cor-2-40}
Equivalent covering families define the same sheaf condition.
\end{proposition}
\begin{definition}\label{fga-def-2-47}\dcref{2.47}\source{fga-explained:page-0085}\uses{fga-def-2-41}
Two Grothendieck topologies are compared by inclusion of their covering sieves;
the finer topology has at least as many covering sieves.
\end{definition}
\begin{proposition}\label{fga-prop-2-48}\dcref{2.48}\source{fga-explained:page-0085}\uses{fga-def-2-47,fga-def-2-37}
If $T\subseteq T'$ then every $T'$-sheaf is a $T$-sheaf; the converse holds
exactly when the topologies have the same covering sieves.
\end{proposition}
\begin{definition}\label{fga-def-2-52}\dcref{2.52}\source{fga-explained:page-0087}\uses{fga-def-2-41}
A topology is saturated when a family is covering whenever its generated sieve
contains a covering sieve.
\end{definition}
\begin{proposition}\label{fga-prop-2-53}\dcref{2.53}\source{fga-explained:page-0087}\uses{fga-def-2-52,fga-cor-2-40}
Saturation does not change the sheaves, and every Grothendieck topology has a
canonical saturated refinement with the same sheaf category.
\end{proposition}
\begin{theorem}\label{fga-thm-2-55}\dcref{2.55}\source{fga-explained:page-0088}\uses{fga-def-2-37,fga-def-2-34,fga-thm-2-1}
Every representable functor on $(\operatorname{Sch}/S)$ is a sheaf for the
Zariski, etale, fppf and fpqc topologies.
\end{theorem}
\begin{proof}
For an fpqc cover $\{U_i\to U\}$, morphisms $U\to X$ are determined by their
pullbacks because the cover is faithfully flat.  Compatible morphisms glue on
affine charts by the equalizer criterion for rings; the glued maps agree on
overlaps and hence glue globally.  The other three topologies are finer
subcases of this argument.
\end{proof}

\begin{definition}\label{fga-def-2-57}\dcref{2.57}\source{fga-explained:page-0095}\uses{fga-thm-2-55}
A topology is subcanonical if every representable functor is a sheaf.
\end{definition}
\begin{definition}\label{fga-def-2-58}\dcref{2.58}\source{fga-explained:page-0095}\uses{fga-def-2-24}
The comma site $(\mathcal C/U)$ has coverings obtained by pullback of coverings
of $U$.
\end{definition}
\begin{proposition}\label{fga-prop-2-59}\dcref{2.59}\source{fga-explained:page-0096}\uses{fga-def-2-57,fga-def-2-58}
On a subcanonical site, the representable functor of an object over $U$ is a
sheaf on the comma site $(\mathcal C/U)$.
\end{proposition}
\begin{lemma}\label{fga-lem-2-60}\dcref{2.60}\source{fga-explained:page-0096}\uses{fga-cor-2-40,fga-def-2-34}
For a scheme $S$, a functor $F:(\operatorname{Sch}/S)^{\mathrm{op}}\to\operatorname{Set}$
is an fpqc sheaf if and only if it satisfies the affine equalizer condition and
is compatible with restriction to affine open covers.
\end{lemma}
\begin{lemma}\label{fga-lem-2-62}\dcref{2.62}\source{fga-explained:page-0098}\uses{fga-prop-1-15}
If $X_1\to Y$ and $X_2\to Y$ are scheme morphisms and their base changes agree
on a faithfully flat cover of $Y$, then they agree over $Y$.
\end{lemma}
\begin{definition}[sheafification]\label{fga-def-2-63}\dcref{2.63}\source{fga-explained:page-0099}\uses{fga-def-2-37}
The sheafification of a presheaf $F$ is a sheaf $F^\#$ equipped with a map
$F\to F^\#$ universal among maps from $F$ to sheaves.
\end{definition}
\begin{theorem}\label{fga-thm-2-64}\dcref{2.64}\source{fga-explained:page-0099}\uses{fga-def-2-63}
Every presheaf on a site has a sheafification, unique up to unique isomorphism;
the construction preserves finite limits needed for group-valued sheaves.
\end{theorem}
\begin{proof}
Iterate the separated-presheaf construction and the plus construction.  The
second plus is a sheaf, and the universal property follows from the equalizer
description of compatible sections.
\end{proof}
